% Clase 21 — curva de primera imantación del hierro inicialmente desmagnetizado,
% como la dibuja el docente: B(H) arranca casi lineal, crece más rápido y
% satura (del orden de 1 T); mu/mu0 tiene un máximo (~6000) cerca de 200 A/m y
% luego decrece. CUALITATIVA: las curvas son funciones de forma, no datos.
% B = 1.75 tanh((H/180)^1.6) T, y mu/mu0 = B/(mu0 H), con máximo ~5900.
\begin{tikzpicture}
  \begin{axis}[
    emfig, width=6.4cm, height=5cm,
    xmin=0, xmax=600, ymin=0, ymax=2.1,
    xtick={0,200,400}, ytick={0,0.5,1,1.5,2},
    xlabel={$H$ (A/m)}, ylabel={$B$ (T)},
    title={\small (a) $B(H)$},
  ]
    \addplot[figblue, domain=0:600, samples=120] {1.75*tanh((x/180)^1.6)};
  \end{axis}
  \begin{axis}[
    emfig, width=6.4cm, height=5cm, xshift=7.2cm,
    xmin=0, xmax=600, ymin=0, ymax=7000,
    xtick={0,200,400}, ytick={0,2000,4000,6000},
    scaled y ticks=false,
    yticklabel style={/pgf/number format/1000 sep={\,}},
    xlabel={$H$ (A/m)}, ylabel={$\mu/\mu_0$},
    title={\small (b) $\mu(H)/\mu_0$},
  ]
    \addplot[figred, domain=1:600, samples=150]
      {1.75*tanh((x/180)^1.6)/(4*pi*1e-7*x)};
  \end{axis}
\end{tikzpicture}
